% ============================================================
% Tabelle zu §3.1 — in Overleaf ablegen als:
%   tables/chap3/chap3.1table.tex
% Aufruf im Fließtext:
%   % ============================================================
% Tabelle zu §3.1 — in Overleaf ablegen als:
%   tables/chap3/chap3.1table.tex
% Aufruf im Fließtext:
%   % ============================================================
% Tabelle zu §3.1 — in Overleaf ablegen als:
%   tables/chap3/chap3.1table.tex
% Aufruf im Fließtext:
%   % ============================================================
% Tabelle zu §3.1 — in Overleaf ablegen als:
%   tables/chap3/chap3.1table.tex
% Aufruf im Fließtext:
%   \input{tables/chap3/chap3.1table}
% Label: t:methodischer-zugang-teilfragen
% Inhalt = Blaupause-Kap-3 §5 Absatz 4 (TF-Zuordnung), sprachlich an
% die tatsächliche Evidenzlage angepasst (keine erfundenen Materialien).
% ============================================================

\begin{table}[tbp]
\centering

\caption[Methodische Zuordnung der Teilfragen]
{Zuordnung der Teilfragen zu Erkenntniszielen, Material und methodischem
Zugang. Eigene Darstellung.}
\label{t:methodischer-zugang-teilfragen}

\small
\renewcommand{\arraystretch}{1.15}

\begin{tabularx}{\textwidth}{
    >{\raggedright\arraybackslash}p{0.09\textwidth}
    >{\raggedright\arraybackslash}p{0.26\textwidth}
    >{\raggedright\arraybackslash}X
}
\toprule
\textbf{Teilfrage} &
\textbf{Erkenntnisziel} &
\textbf{Material und methodischer Zugang} \\
\midrule

TF1 &
Herausforderungen und Designanforderungen &
Nutzungsziel, ausgewählte frühe Entwicklungsstände, Läufe und
Entscheidungsnotizen;
formative Rekonstruktion und Design Rationale (Kapitel~\ref{chap:explorative-problemklaerung}) \\
\addlinespace[0.4em]

TF2 &
Gestaltung und Realisierung &
finaler Artefaktstand, Architektur und Frameworkabbildung;
Artefaktbeschreibung und technische Demonstration (Kapitel~\ref{chap:loesungsarchitektur} und~\ref{chap:maf-realisierung}) \\
\addlinespace[0.4em]

TF3 &
Eigenschaften und Trade-offs &
eingefrorene Evaluationsläufe, Referenzannotation und Challenge-Tests;
quantitative und qualitative Evaluation (Kapitel~\ref{chap:evaluation}) \\

\bottomrule
\end{tabularx}
\end{table}

% Label: t:methodischer-zugang-teilfragen
% Inhalt = Blaupause-Kap-3 §5 Absatz 4 (TF-Zuordnung), sprachlich an
% die tatsächliche Evidenzlage angepasst (keine erfundenen Materialien).
% ============================================================

\begin{table}[tbp]
\centering

\caption[Methodische Zuordnung der Teilfragen]
{Zuordnung der Teilfragen zu Erkenntniszielen, Material und methodischem
Zugang. Eigene Darstellung.}
\label{t:methodischer-zugang-teilfragen}

\small
\renewcommand{\arraystretch}{1.15}

\begin{tabularx}{\textwidth}{
    >{\raggedright\arraybackslash}p{0.09\textwidth}
    >{\raggedright\arraybackslash}p{0.26\textwidth}
    >{\raggedright\arraybackslash}X
}
\toprule
\textbf{Teilfrage} &
\textbf{Erkenntnisziel} &
\textbf{Material und methodischer Zugang} \\
\midrule

TF1 &
Herausforderungen und Designanforderungen &
Nutzungsziel, ausgewählte frühe Entwicklungsstände, Läufe und
Entscheidungsnotizen;
formative Rekonstruktion und Design Rationale (Kapitel~\ref{chap:explorative-problemklaerung}) \\
\addlinespace[0.4em]

TF2 &
Gestaltung und Realisierung &
finaler Artefaktstand, Architektur und Frameworkabbildung;
Artefaktbeschreibung und technische Demonstration (Kapitel~\ref{chap:loesungsarchitektur} und~\ref{chap:maf-realisierung}) \\
\addlinespace[0.4em]

TF3 &
Eigenschaften und Trade-offs &
eingefrorene Evaluationsläufe, Referenzannotation und Challenge-Tests;
quantitative und qualitative Evaluation (Kapitel~\ref{chap:evaluation}) \\

\bottomrule
\end{tabularx}
\end{table}

% Label: t:methodischer-zugang-teilfragen
% Inhalt = Blaupause-Kap-3 §5 Absatz 4 (TF-Zuordnung), sprachlich an
% die tatsächliche Evidenzlage angepasst (keine erfundenen Materialien).
% ============================================================

\begin{table}[tbp]
\centering

\caption[Methodische Zuordnung der Teilfragen]
{Zuordnung der Teilfragen zu Erkenntniszielen, Material und methodischem
Zugang. Eigene Darstellung.}
\label{t:methodischer-zugang-teilfragen}

\small
\renewcommand{\arraystretch}{1.15}

\begin{tabularx}{\textwidth}{
    >{\raggedright\arraybackslash}p{0.09\textwidth}
    >{\raggedright\arraybackslash}p{0.26\textwidth}
    >{\raggedright\arraybackslash}X
}
\toprule
\textbf{Teilfrage} &
\textbf{Erkenntnisziel} &
\textbf{Material und methodischer Zugang} \\
\midrule

TF1 &
Herausforderungen und Designanforderungen &
Nutzungsziel, ausgewählte frühe Entwicklungsstände, Läufe und
Entscheidungsnotizen;
formative Rekonstruktion und Design Rationale (Kapitel~\ref{chap:explorative-problemklaerung}) \\
\addlinespace[0.4em]

TF2 &
Gestaltung und Realisierung &
finaler Artefaktstand, Architektur und Frameworkabbildung;
Artefaktbeschreibung und technische Demonstration (Kapitel~\ref{chap:loesungsarchitektur} und~\ref{chap:maf-realisierung}) \\
\addlinespace[0.4em]

TF3 &
Eigenschaften und Trade-offs &
eingefrorene Evaluationsläufe, Referenzannotation und Challenge-Tests;
quantitative und qualitative Evaluation (Kapitel~\ref{chap:evaluation}) \\

\bottomrule
\end{tabularx}
\end{table}

% Label: t:methodischer-zugang-teilfragen
% Inhalt = Blaupause-Kap-3 §5 Absatz 4 (TF-Zuordnung), sprachlich an
% die tatsächliche Evidenzlage angepasst (keine erfundenen Materialien).
% ============================================================

\begin{table}[tbp]
\centering

\caption[Methodische Zuordnung der Teilfragen]
{Zuordnung der Teilfragen zu Erkenntniszielen, Material und methodischem
Zugang. Eigene Darstellung.}
\label{t:methodischer-zugang-teilfragen}

\small
\renewcommand{\arraystretch}{1.15}

\begin{tabularx}{\textwidth}{
    >{\raggedright\arraybackslash}p{0.09\textwidth}
    >{\raggedright\arraybackslash}p{0.26\textwidth}
    >{\raggedright\arraybackslash}X
}
\toprule
\textbf{Teilfrage} &
\textbf{Erkenntnisziel} &
\textbf{Material und methodischer Zugang} \\
\midrule

TF1 &
Herausforderungen und Designanforderungen &
Nutzungsziel, ausgewählte frühe Entwicklungsstände, Läufe und
Entscheidungsnotizen;
formative Rekonstruktion und Design Rationale (Kapitel~\ref{chap:explorative-problemklaerung}) \\
\addlinespace[0.4em]

TF2 &
Gestaltung und Realisierung &
finaler Artefaktstand, Architektur und Frameworkabbildung;
Artefaktbeschreibung und technische Demonstration (Kapitel~\ref{chap:loesungsarchitektur} und~\ref{chap:maf-realisierung}) \\
\addlinespace[0.4em]

TF3 &
Eigenschaften und Trade-offs &
eingefrorene Evaluationsläufe, Referenzannotation und Challenge-Tests;
quantitative und qualitative Evaluation (Kapitel~\ref{chap:evaluation}) \\

\bottomrule
\end{tabularx}
\end{table}
